\section{Variáveis Aleatórias Contínuas}

Variáveis aleatória discretas caracterizam-se por uma lista enumerável de pontos
concentrarem massas de probabilidades.
Outrossim, variáveis aleatórias contínuas são aquelas em que nenhum ponto real
acumula qualquer massa de probabilidade, estando essa distribuída ao longo de suas
vizinhanças.
Formalmente, teremos que qualquer conjunto de pontos de medida nula (como pontos contáveis)
terá probabilidade desprezível, caracterizando a medida de probablidade da variável como
absolutamente contínua em relação a medida de Lebesgue
(ver Definição \ref{def:medida absolutamente continua}).

\begin{def:variavel continua}
A variável aleatória $X$ será contínua se $P_X \ll \lambda$.
\end{def:variavel continua}

O Teorema \nameref{thm:radon nikodym} garante a existência de uma função especial
-- mensurável, não negativa e Lebesgue-integrável -- associada a medida de probablidade
induzida $P_X$ de $X$.

\begin{prop:existencia fdp}
Toda variável aleatória contínua $X$ admite uma função real mensurável não negativa
$f \in L^1(\lambda)$ tal que
\[
	f = \frac{d}{d\lambda} P_X
\]
\begin{proof}
	Se $X$ é contínua, então $P_X \ll \lambda$. Como $\lambda$ é $\sigma$-finita, então
	pelo Teorema \nameref{thm:radon nikodym} existe $f \in L^1(\lambda)$ não negativa
	tal que
	\[
		f = \frac{d}{d\lambda} P_X
	\]
\end{proof}
\end{prop:existencia fdp}

Se interpretarmos a probabilidade como uma massa distribuída em espaço (neste caso,
em $\real$) e que quanto maior a extensão do intervalo avaliado, maior a probabilidade
da variável assumir realizações nele, então podemos interpretar a função $f$ como
uma \textbf{densidade} de probabilidade.

\begin{def:fdp}
Seja $X$ uma variável aleatória contínua com medida induzida $P_X$. A função $f$
que satisfaz
\[
	f = \frac{d}{d\lambda}P_X \Longleftrightarrow P_X = \int f\, d\lambda
\]
é chamada de densidade de probabilidade de $X$.
\end{def:fdp}

A função de densidade pode ser então usada para medir a probabilidade de uma variável
contínua assumir realizações num intervalo de valores reais.
Isso é feito avaliando a integral da função sobre o intervalo desejado.
A relação mais direta entre as funções de distribuição e de densidade é dada pela
Proposição a seguir.

\begin{prop:fdp e fda}
\label{prop:fdp e fda}
Se $X$ é uma variável contínua com densidade $f$ e distribuição $F$, então
\[
	F(x) = \int_{-\infty}^x f(t)\,dt
\]
\begin{proof}
	\[
		F(x) = P_X((-\infty, x]) = \int_{(-\infty, x]} f\, d\lambda = \int_{-\infty}^x f(t)\, dt
	\]
\end{proof}
\end{prop:fdp e fda}

A distribuições de variáveis discretas tem a forma de uma função constante por partes,
pois entre um ponto e outro temos um ``salto'' que representa a probabilidade acumulada
por um dos pontos.
Em contraste, como para variáveis contínuas nenhum ponto acumula massa de probabilidade,
a sua distrbiuição deve variar suavemente ao longo de qualquer intervalo, ou de maneira mais
forte, em qualquer conjunto disjunto de intervalso
Isso caracteriza a distribuição como absolutamente contínua (ver Definição
\ref{def:absolutamente continua}).

\begin{prop:continuidade fda}
A variável $X$ é contínua se, e somente se, sua distribuição $F$ é absolutamente contínua.
\end{prop:continuidade fda}
\begin{proof}
	Supondo que $X$ é contínua, então $P_X \ll \lambda$.
	Pela Proposição \ref{prop:steltjes e fda} e pelo Teorema \ref{thm:medida steltjes},
	concluiremos que $F$ é absolutamente contínua.
	Por outro lado, se $F$ é absolutamente contínua, então $\mu_F = P_X \ll \lambda$ e,
	por definição, a variável $X$ será contínua.
\end{proof}

Como qualquer conjunto de pontos isolados terá probabilidade nula, então é intuitivo
pensar que o intervalo $(a, b]$ terá a mesma probabilidade do que $(a, b)$.
A formalização dessas ideias é dada pelo Teorema a seguir.

\begin{thm:prob continua}
\label{thm:prob continua}
Se $X$ é uma variável aleatória contínua com densidade $f$,
então para todo $x, a, b \in \real$
\begin{enumerate}
	\item[i.] $P(a < X \leq b) = \int_a^b f(t)\,dt$
	\item[ii.] $P(X = x) = 0$
	\item[iii.] $P(X \leq x) = P(X < x)$
	\item[iv.] $P(X \geq x) = P(X > x)$
\end{enumerate}
\begin{proof}
	\[
		P(a < X \le b) = P(X \le b) - P(X \le a) = F(b) - F(a)
	\]
	e pela Proposição \ref{prop:fdp e fda}
	\[
		P(a < X \le b) = \int_{-\infty}^b f(t)\,dt - \int_{-\infty}^a f(t)\,dt = \int_a^b f(t)\,dt
	\]
	provando (i).

	Observemos agora que
	\[
		\lambda(\{x\}) = 0 \Longrightarrow \int_{\{x\}} f\, d\lambda = 0 \Longrightarrow P(X = x) = 0
	\]
	provando (ii).

	Uma vez que
	\[
		P(X \leq x) = P(X = x) + P(X < x) \quad \text{e} \quad P(X \ge x) = P(X = x) + P(X > x)
	\]
	então (iii) e (iv) são consequências de (ii).
\end{proof}
\end{thm:prob continua}

\begin{prop:densidade 1}
\label{prop:densidade 1}
Se $X$ é uma variável aleatória contínua com densidade $f$ e suporte $\Omega_X$, então
\[
	\int_{\Omega_X} f\,d\lambda = \int_\real f\,\lambda = 1
\]
\begin{proof}
	Com efeito,
	\[
		P_X(\Omega_X) = 1 \Longrightarrow \int_{\Omega_X}f\, d\lambda = 1.
	\]
	Ademais, se $(\real, \borel, P_X)$ é o espaço de probabilidade induzido por
	$X$, então
	\[
		P_X(\real) = P(\Omega) = 1 \Longrightarrow \int_\real f\, d\lambda = 1
	\]
\end{proof}
\end{prop:densidade 1}

\begin{prop:densidade fora suporte}
Se $X$ é uma variável contínua com densidade $f$ e suporte $\Omega_X$, então
\[
	f = 0 \qs \text{ em } \real \setminus \Omega_X
\]
\begin{proof}
	Seja $A = \real \setminus \Omega_X$. Observe que
	\[
		\int_{A} f\, d\lambda + \int_{\Omega_X} f\, d\lambda = \int_\real f\, d\lambda
	\]
	e pela Proposição \ref{prop:densidade 1}, conclui-se
	\[
		\int_A f\, d\lambda = 0.
	\]
	Como um função de densidade é mensurável e não negativa, então por nulidade
	da integral de Lebesgue, obtém-se que $f = 0 \qs$ em $A$.
\end{proof}
\end{prop:densidade fora suporte}

\begin{lem:densidade continua}
\label{lem:densidade continua}
Seja $X$ variável contínua com densidade $f$ e suporte $\Omega_X$.
Se $f$ é contínua em $x \in \real$, então $x \in \Omega_X$.
\begin{proof}
	Suponha que $0 < \epsilon < f(x)$. Por continuidade em $x$, existe $\delta > 0$ tal que
	\[
		\forall a \in \real, |x-a| < \delta \Longrightarrow |f(x)-f(a)| < \epsilon
	\]
	Suponha por contradição que $x \notin \Omega_X$.
	Com efeito, exista uma vizinhança $V = (x-\nu, x+\nu)$, com $\nu > 0$
	tal que $P(X \in V) = 0$, que equivale a $f=0 \qs$ em $V$.
	Existe, pois, $y \in V$ tal que
	\[
		|x-y| < \delta \quad \text{e} \quad f(y) = 0
	\]
	de modo que
	\[
		|f(x) - f(y)| < \epsilon \Longrightarrow |f(x)| < \epsilon,
	\]
	e sendo $f$ densidade, $f > 0$ e $|f(x)| = f(x)$.
	Por conseguinte,
	\[
		f(x) < \epsilon
	\]
	contradizendo a afirmação que $\epsilon < f(x)$.
	Conclui-se que $x \in \Omega_X$.
\end{proof}
\end{lem:densidade continua}

\begin{thm:suporte continua}
\label{thm:suporte continua}
Se $X$ é variável aleatória com densidade $f$ contínua em
\[
	\{x \in \real \mid f(x) > 0\}
\]
então
\[
	\Omega_X = \overline{\{x \in \real \mid f(x) > 0\}}
\]
\begin{proof}
	Seja
	\[
		E = \{x \in \real \mid f(x) > 0\}.
	\]
	Mostremos que $\Omega_X \subset \overline E$.
	Para tanto, suponha que $x \notin \overline E$.
	Existe, então, vizinhança $V = (x-\delta, x+\delta)$, com $\delta > 0$, tal que
	\[
		V \cap E = \varnothing
	\]
	Então para todo $a \in V$, $f(a) = 0$, o que equivale a dizer que $f = 0 \qs$ em $V$.
	Por consequência, $P(X \in V) = 0$, e $x \notin \Omega_X$.

	Demonstremos agora que $\overline E \subset \Omega_X$.
	Do Lema \ref{lem:densidade continua}, temos que $E \subset \Omega_X$.
	Resta mostrar que todo ponto de acumulação $x \notin E$ é membro de $\Omega_X$.
	Se $x \in \overline E$, então
	\[
		\forall \nu > 0, (x-\nu, x+\nu) \cap E \ne \varnothing.
	\]
	Isso equivale a dizer que para toda vizinhança $V$ de $x$ contém pelo menos
	um ponto com densidade positiva e, por consequência do Lema \ref{lem:densidade continua},
	que é membro do suporte.
	Se tivéssemos $x \notin \Omega_X$, então haveria uma vizinhança $U$
	centrada em $x$ tal que $P(X \in U) = 0$, e teríamos
	\[
		\forall a \in U, a \notin \Omega_X
	\]
	contradizendo a afirmação que toda vizinhança de $x$ tem um membro do suporte.
	Portanto, $\overline E \subset \Omega_X$.
\end{proof}
\end{thm:suporte continua}
